\chapter{Sobolev Spaces}

\section{One-dimensional Sobolev spaces}

\begin{notedefinition}
Let $I\subseteq\mathbb R$ be an open interval and let $1\le p\le\infty$.  A
function $u\in L^p(I)$ belongs to $W^{1,p}(I)$ if there is $v\in L^p(I)$ such
that
\[
  \int_I u\varphi'\,\dd x
  =-\int_I v\varphi\,\dd x
  \qquad(\varphi\in C_c^\infty(I)).
\]
The function $v$ is unique almost everywhere and is denoted by $u'$.  The norm
is
\[
  \|u\|_{W^{1,p}(I)}=\|u\|_{L^p(I)}+\|u'\|_{L^p(I)}.
\]
We write $H^1(I)=W^{1,2}(I)$.
\end{notedefinition}

\begin{notetheorem}
The space $W^{1,p}(I)$ is Banach for every $1\le p\le\infty$, separable for
$1\le p<\infty$, and reflexive for $1<p<\infty$.
\end{notetheorem}

\begin{proof}
The map
\[
  u\longmapsto(u,u')
\]
is an isometry from $W^{1,p}(I)$ into $L^p(I)\times L^p(I)$.  Its range is
closed: if $u_n\to u$ and $u_n'\to v$ in $L^p$, then for every test function
\[
  \int u\varphi'
  =\lim_n\int u_n\varphi'
  =-\lim_n\int u_n'\varphi
  =-\int v\varphi.
\]
Completeness, separability, and reflexivity therefore follow from the
corresponding product-space properties of $L^p$.
\end{proof}
\begin{notetheorem}[One-dimensional fundamental theorem]
Let $u\in W^{1,p}(I)$.  There is an absolutely continuous representative
$\widetilde u$ such that
\[
  \widetilde u(x)-\widetilde u(y)
  =\int_y^x u'(t)\,\dd t
  \qquad(x,y\in I).
\]
In particular, every one-dimensional Sobolev function has a continuous
representative.
\end{notetheorem}

\begin{proof}
Fix $y_0\in I$ and define
\[
  v(x)=\int_{y_0}^x u'(t)\,\dd t.
\]
Then $v$ is absolutely continuous and has weak derivative $u'$.  Hence
$u-v$ has weak derivative zero.  The lemma below shows that $u-v$ is almost
everywhere constant, so $u$ agrees almost everywhere with $v+c$.
\end{proof}

\begin{notelemma}
If $f\in L^1_{\mathrm{loc}}(I)$ and
\[
  \int_I f\varphi'\,\dd x=0
  \qquad(\varphi\in C_c^\infty(I)),
\]
then $f$ is almost everywhere constant on $I$.
\end{notelemma}

\begin{proof}
Choose $\psi\in C_c^\infty(I)$ with $\int_I\psi=1$.  For any
$\varphi\in C_c^\infty(I)$, the function
$\varphi-(\int\varphi)\psi$ has integral zero and hence is the derivative of
a compactly supported smooth function.  Therefore
\[
  \int f\varphi
  =\left(\int f\psi\right)\left(\int\varphi\right),
\]
which says that the distribution defined by $f$ is the constant
$c=\int f\psi$.
\end{proof}
\begin{noteproposition}[Dual characterization of a weak derivative]
Let $u\in L^p(I)$ and $1<p\le\infty$.  Then $u\in W^{1,p}(I)$ if and only
if there is $C>0$ such that
\[
  \left|\int_Iu\varphi'\,\dd x\right|
  \le C\|\varphi\|_{L^{p'}(I)}
  \qquad(\varphi\in C_c^\infty(I)),
\]
with $p'=1$ when $p=\infty$.  The smallest such $C$ is
$\|u'\|_p$.
\end{noteproposition}

\begin{proof}
One direction is Hölder's inequality.  Conversely, the map
$\varphi\mapsto-\int u\varphi'$ is bounded on the dense subspace
$C_c^\infty(I)\subset L^{p'}(I)$.  Extend it to $L^{p'}$ and apply the
$L^p$ duality theorem to obtain $v\in L^p$ with
$-\int u\varphi'=\int v\varphi$.
\end{proof}

\begin{noteremark}
At $p=1$, the displayed estimate with the $L^\infty$ norm controls the
distributional derivative only as a finite measure and characterizes bounded
variation.  It does not by itself force that measure to have an $L^1$
density.  This is why the theorem is stated for $p>1$ (and for $p=\infty$ via
$(L^1)^*=L^\infty$).
\end{noteremark}
\begin{noteremark}
For $p=1$, the absolutely continuous representative is exactly the classical
one: a function belongs to $W^{1,1}(I)$ precisely when it is almost everywhere
equal to an absolutely continuous function whose classical derivative lies in
$L^1$.  Merely bounded variation is weaker; BV functions may have jump or
singular derivatives.
\end{noteremark}

\begin{notetheorem}[Lipschitz characterization]
A function $u\in L^\infty(I)$ belongs to $W^{1,\infty}(I)$ if and only if it
has a Lipschitz representative.  Moreover,
\[
  \operatorname{Lip}(u)=\|u'\|_\infty.
\]
\end{notetheorem}

\begin{proof}
The fundamental theorem gives
$|u(x)-u(y)|\le\|u'\|_\infty|x-y|$.  Conversely, a Lipschitz function is
absolutely continuous, differentiable almost everywhere, and its derivative
is bounded by its Lipschitz constant.  Integration by parts then identifies
that derivative as the weak derivative.
\end{proof}
\begin{notetheorem}[Difference-quotient criterion on $\mathbb R$]
Let $1<p<\infty$ and $u\in L^p(\mathbb R)$.  Then $u\in W^{1,p}(\mathbb R)$
if and only if there is $C>0$ such that
\[
  \|u(\cdot+h)-u\|_p\le C|h|
  \qquad(h\in\mathbb R).
\]
The least admissible $C$ equals $\|u'\|_p$.
\end{notetheorem}

\begin{proof}
If $u\in W^{1,p}$, then for the absolutely continuous representative,
\[
  u(x+h)-u(x)=\int_0^h u'(x+s)\,\dd s.
\]
Minkowski's integral inequality gives the estimate with $C=\|u'\|_p$.
Conversely, for $\varphi\in C_c^\infty(\mathbb R)$,
\[
  \int u(x)\frac{\varphi(x-h)-\varphi(x)}{h}\,\dd x
  =\int\frac{u(x+h)-u(x)}{h}\varphi(x)\,\dd x.
\]
The right-hand side is bounded by $C\|\varphi\|_{p'}$.  Letting $h\to0$ and
using the preceding dual characterization yields a weak derivative in $L^p$.
For $p=1$, the analogous uniform translation bound characterizes bounded
variation rather than $W^{1,1}$, so the endpoint cannot be inserted without
an additional absolute-continuity condition.
\end{proof}
\section{Extension and smooth approximation}

\begin{notetheorem}[Extension from an interval]
Let $I$ be an interval and $1\le p\le\infty$.  There is a bounded linear
operator
\[
  E\colon W^{1,p}(I)\to W^{1,p}(\mathbb R)
\]
such that $Eu=u$ almost everywhere on $I$ and
\[
  \|Eu\|_{W^{1,p}(\mathbb R)}
  \le C_I\|u\|_{W^{1,p}(I)}.
\]
For a bounded interval, $Eu$ may be chosen with support in a fixed bounded
neighbourhood of $I$.
\end{notetheorem}

\begin{proof}[Construction]
For a half-line, reflect across the endpoint on a short adjacent interval and
then multiply by a fixed smooth cutoff.  Reflection preserves the $L^p$ norms
of both the function and its weak derivative up to a fixed factor.  For a
bounded interval, reflect across both endpoints, use a partition of unity near
the two ends, and multiply by a cutoff supported in a slightly larger
interval.  The product rule gives the stated bound.
\end{proof}
\begin{notelemma}[Mollification of Sobolev functions]
Let $u\in W^{1,p}(\mathbb R)$ and let $\rho\in C_c^\infty(\mathbb R)$.  Then
\[
  \rho*u\in W^{1,p}(\mathbb R),
  \qquad
  (\rho*u)'=\rho*u'.
\]
\end{notelemma}

\begin{proof}
For a test function $\varphi$, Fubini's theorem and the definition of the weak
derivative give
\[
 \int(\rho*u)\varphi'
 =\int u\,(\widetilde\rho*\varphi')
 =-\int u'\,(\widetilde\rho*\varphi)
 =-\int(\rho*u')\varphi,
\]
where $\widetilde\rho(x)=\rho(-x)$.  Young's inequality gives the $L^p$
bounds.
\end{proof}

\begin{notetheorem}[Smooth approximation]
For $1\le p<\infty$, $C_c^\infty(\mathbb R)$ is dense in
$W^{1,p}(\mathbb R)$.  More generally, smooth functions are dense in
$W^{1,p}(I)$, and on a bounded interval they may be obtained by extending,
mollifying, and restricting.
\end{notetheorem}

\begin{proof}
Let $\rho_n$ be mollifiers.  Then
\[
  \rho_n*u\to u,
  \qquad
  \rho_n*u'\to u'
\]
in $L^p$, so $\rho_n*u\to u$ in $W^{1,p}$.  Multiplying by cutoffs that tend
to one produces compact support while the product-rule error tends to zero.
For an interval, first apply the bounded extension operator.
\end{proof}
\section{One-dimensional Sobolev embeddings}

\begin{notetheorem}[Sobolev embedding on an interval]
Let $I$ be an interval of finite length and $1\le p\le\infty$.  Every
$u\in W^{1,p}(I)$ has a continuous representative and
\[
  \|u\|_{L^\infty(I)}
  \le C_I\|u\|_{W^{1,p}(I)}.
\]
Thus $W^{1,p}(I)\hookrightarrow L^\infty(I)$ continuously.

If $I$ is bounded, then
\begin{enumerate}
  \item $W^{1,p}(I)\hookrightarrow C(\overline I)$ compactly for
        $1<p\le\infty$;
  \item $W^{1,1}(I)\hookrightarrow L^2(I)$ compactly.
\end{enumerate}
\end{notetheorem}

\begin{proof}
For $1<p<\infty$, apply the fundamental theorem to
$G(u)=|u|^{p-1}u$ (or a smooth approximation in the complex case):
\[
  |u(x)|^p
  \le p\int_I|u|^{p-1}|u'|+\frac1{|I|}\|u\|_p^p,
\]
and Hölder's inequality gives the stated $L^\infty$ estimate.  The endpoint
$p=1$ follows directly from
$|u(x)|\le |I|^{-1}\|u\|_1+\|u'\|_1$, and $p=\infty$ is immediate.

For $p>1$, a bounded set in $W^{1,p}(I)$ is uniformly bounded and satisfies
\[
  |u(x)-u(y)|
  \le\|u'\|_p|x-y|^{1/p'},
\]
so Arzelà--Ascoli gives compactness in $C(\overline I)$.  For $p=1$, use the
extension operator.  A bounded family is bounded in both $L^1$ and
$L^\infty$, and translations satisfy
$\|u(\cdot+h)-u\|_1\le C|h|$.  The Fréchet--Kolmogorov compactness criterion
gives relative compactness in $L^1$; interpolation with the uniform
$L^\infty$ bound yields relative compactness in $L^2$.
\end{proof}
\begin{notecorollary}[Decay on an unbounded interval]
Let $I$ be unbounded and let $1\le p<\infty$.  If $u\in W^{1,p}(I)$, then its
continuous representative satisfies
\[
  u(x)\longrightarrow0
\]
at every infinite endpoint of $I$.
\end{notecorollary}

\begin{proof}
Approximate $u$ in $W^{1,p}(I)$ by compactly supported smooth functions.  The
Sobolev $L^\infty$ estimate makes this approximation uniform, so $u$ is a
uniform limit of functions vanishing at infinity.
\end{proof}

\begin{noteremark}
On an unbounded interval, the embedding
$W^{1,p}(I)\hookrightarrow L^\infty(I)$ is not compact.  Translates of one
fixed nonzero compactly supported smooth function are bounded in $W^{1,p}$ but
have no uniformly convergent subsequence.
\end{noteremark}
\section{Product and chain rules}

\begin{noteproposition}[One-dimensional differentiation rules]
Let $1\le p\le\infty$.
\begin{enumerate}
  \item If $u,v\in W^{1,p}(I)\cap L^\infty(I)$, then
        $uv\in W^{1,p}(I)$ and
        \[
          (uv)'=u'v+uv'.
        \]
        The integration-by-parts identity
        \[
          \int_a^b u'v
          =u(b)v(b)-u(a)v(a)-\int_a^buv'
        \]
        holds for bounded $I=(a,b)$.
  \item If $G\in C^1(\mathbb R)$, $G(0)=0$, and $G'$ is bounded, then
        $G\circ u\in W^{1,p}(I)$ and
        \[
          (G\circ u)'=(G'\circ u)u'.
        \]
\end{enumerate}
\end{noteproposition}

\begin{proof}
Use smooth approximation in $W^{1,p}$ together with the one-dimensional
$L^\infty$ embedding on bounded subintervals.  The classical identities hold
for the approximants, and the products converge in $L^p$.  For the chain rule,
first handle bounded $G'$ directly and pass to the limit; local truncation
handles the corresponding local version.
\end{proof}
\section{Higher-order spaces, traces, and negative norms}

\begin{notedefinition}
For an integer $m\ge1$,
\[
  W^{m,p}(I)
  =\{u\in L^p(I):u^{(j)}\in L^p(I),\ 1\le j\le m\},
\]
with norm
\[
  \|u\|_{W^{m,p}(I)}
  =\sum_{j=0}^m\|u^{(j)}\|_p.
\]
Equivalent $\ell^p$ combinations of these terms give the same topology.
When $p=2$ we write $H^m(I)=W^{m,2}(I)$.
\end{notedefinition}

\begin{notedefinition}
For $1\le p<\infty$,
\[
  W_0^{1,p}(I)=\overline{C_c^\infty(I)}^{\,W^{1,p}(I)}.
\]
\end{notedefinition}

\begin{notetheorem}[Trace characterization]
Let $I=(a,b)$ be bounded.  Then
\[
  W_0^{1,p}(I)
  =\{u\in W^{1,p}(I):u(a)=u(b)=0\},
\]
where the boundary values are taken from the continuous representative.
\end{notetheorem}

\begin{proof}
Trace continuity follows from the Sobolev embedding, so every limit of compactly
supported smooth functions has zero endpoint values.  Conversely, if $u$ has
zero endpoints, multiply it by cutoffs that equal one away from successively
smaller endpoint intervals.  The zero trace and the fundamental theorem make
the cutoff errors tend to zero in $W^{1,p}$; mollification then gives
$C_c^\infty$ approximation.
\end{proof}
\begin{notetheorem}[Poincaré inequality]
For a bounded interval $I$ and $1\le p<\infty$, there is $C_I>0$ such that
\[
  \|u\|_p\le C_I\|u'\|_p
  \qquad(u\in W_0^{1,p}(I)).
\]
The same estimate holds for zero-trace functions in $W^{1,\infty}(I)$.
Consequently, $\|u'\|_p$ is an equivalent norm on $W_0^{1,p}(I)$.
\end{notetheorem}

\begin{proof}
For the continuous representative and $u(a)=0$,
\[
  |u(x)|\le\int_a^x|u'(t)|\,\dd t.
\]
Hölder's inequality and integration in $x$ give the result.
\end{proof}
\begin{notedefinition}
For $m\ge1$, $W_0^{m,p}(I)$ is the closure of $C_c^\infty(I)$ in
$W^{m,p}(I)$.  For $1<p<\infty$, the dual of $W_0^{1,p}(I)$ is denoted
\[
  W^{-1,p'}(I),
\]
and the dual of $H_0^1(I)$ is denoted $H^{-1}(I)$.
\end{notedefinition}

\begin{notetheorem}[Representation of $W^{-1,p'}$]
Let $1<p<\infty$.  A functional $F$ belongs to $W^{-1,p'}(I)$ if and only if
there are $f_0,f_1\in L^{p'}(I)$ such that
\[
  \dual{F}{u}
  =\int_I f_0u\,\dd x+\int_I f_1u'\,\dd x
  \qquad(u\in W_0^{1,p}(I)).
\]
Moreover,
\[
  \|F\|_{W^{-1,p'}}
  =\inf\bigl\{
       \|f_0\|_{p'}+\|f_1\|_{p'}:
       (f_0,f_1)\text{ represents }F
     \bigr\}
\]
up to the harmless choice of an equivalent product norm.  If $I$ is bounded,
Poincaré's inequality allows a representation with $f_0=0$.
\end{notetheorem}

\begin{proof}
Embed $W_0^{1,p}(I)$ isometrically as the closed subspace
\[
  G=\{(u,u'):u\in W_0^{1,p}(I)\}
  \subset L^p(I)\times L^p(I).
\]
The functional $F$ induces a bounded functional on $G$.  Hahn--Banach extends
it to the product, and $L^p$ duality represents that extension by
$(f_0,f_1)\in L^{p'}\times L^{p'}$.  Conversely, Hölder's inequality makes
every such formula continuous.  On a bounded interval, use the derivative
embedding $u\mapsto u'$ and Poincaré's inequality instead.
\end{proof}
\section{Sobolev spaces in several variables}

\begin{notedefinition}
Let $\Omega\subseteq\mathbb R^N$ be open.  A function $u\in L^p(\Omega)$
belongs to $W^{1,p}(\Omega)$ if, for each $1\le i\le N$, there is
$g_i\in L^p(\Omega)$ satisfying
\[
  \int_\Omega u\,\partial_i\varphi\,\dd x
  =-\int_\Omega g_i\varphi\,\dd x
  \qquad(\varphi\in C_c^\infty(\Omega)).
\]
We write $\partial_i u=g_i$ and
$\nabla u=(\partial_1u,\ldots,\partial_Nu)$.  The norm
\[
  \|u\|_{W^{1,p}(\Omega)}
  =\|u\|_p+\sum_{i=1}^N\|\partial_i u\|_p
\]
may be replaced by any equivalent finite-dimensional combination of these
terms.
\end{notedefinition}

\begin{noteproposition}
The space $W^{1,p}(\Omega)$ is Banach for every $1\le p\le\infty$, separable
for $1\le p<\infty$, and reflexive for $1<p<\infty$.  The space
$H^1(\Omega)$ is Hilbert under
\[
  \inner{u}{v}_{H^1}
  =\int_\Omega u\overline v\,\dd x
   +\sum_{i=1}^N\int_\Omega
      \partial_i u\,\overline{\partial_i v}\,\dd x.
\]
\end{noteproposition}

\begin{proof}
Use the closed graph of
$u\mapsto(u,\partial_1u,\ldots,\partial_Nu)$ in a finite product of $L^p$
spaces, exactly as in one dimension.
\end{proof}
\begin{noteproposition}[Closedness of weak differentiation]
Suppose $u_n\to u$ in $L^p(\Omega)$ and
$\partial_i u_n\to g_i$ in $L^p(\Omega)$ for every $i$.  Then
$u\in W^{1,p}(\Omega)$ and $\partial_i u=g_i$.  The same conclusion holds
when all convergences are weak in the corresponding $L^p$ spaces.
\end{noteproposition}

\begin{notedefinition}
For open sets $\omega,\Omega\subseteq\mathbb R^N$, the notation
$\omega\Subset\Omega$ means that $\overline\omega$ is compact and contained
in $\Omega$.
\end{notedefinition}

\begin{notelemma}[Local mollification]
If $u\in W^{1,p}(\Omega)$ and $\omega\Subset\Omega$, then for sufficiently
small $\varepsilon$,
\[
  u_\varepsilon=\rho_\varepsilon*u
\]
is smooth on $\omega$ and
\[
  \partial_i u_\varepsilon
  =\rho_\varepsilon*\partial_i u
  \quad\text{on }\omega.
\]
\end{notelemma}

\begin{notetheorem}[Meyers--Serrin density theorem]
Let $1\le p<\infty$.  Then
\[
  C^\infty(\Omega)\cap W^{1,p}(\Omega)
\]
is dense in $W^{1,p}(\Omega)$.  If $\Omega=\mathbb R^N$, the approximants may
be chosen in $C_c^\infty(\mathbb R^N)$.  On a general open set,
$C_c^\infty(\Omega)$ is dense only in the smaller space
$W_0^{1,p}(\Omega)$, not in all of $W^{1,p}(\Omega)$.
\end{notetheorem}

\begin{proof}[Proof sketch]
Choose an exhaustion $\Omega_1\Subset\Omega_2\Subset\cdots$ and a smooth
partition of unity subordinate to annular pieces.  Mollify each localized
piece at a scale smaller than its distance from the boundary and choose the
scales so that the sum of the $W^{1,p}$ errors is finite and tends to zero.
On $\mathbb R^N$, first cut off at large radius and then mollify.
\end{proof}
\begin{notetheorem}[Distribution and translation criteria]
Let $1<p<\infty$ and $u\in L^p(\mathbb R^N)$.  The following are equivalent:
\begin{enumerate}
  \item $u\in W^{1,p}(\mathbb R^N)$;
  \item for each $i$ there is $C_i$ such that
        \[
          \left|\int_{\mathbb R^N}u\,\partial_i\varphi\,\dd x\right|
          \le C_i\|\varphi\|_{p'}
          \qquad(\varphi\in C_c^\infty);
        \]
  \item there is $C>0$ such that
        \[
          \|u(\cdot+h)-u\|_p\le C|h|
          \qquad(h\in\mathbb R^N).
        \]
\end{enumerate}
For $p=1$, the analogous bounded distributional-derivative or translation
condition characterizes $BV$, not necessarily $W^{1,1}$.
\end{notetheorem}
\begin{noteproposition}[Differentiation rules in $\mathbb R^N$]
Let $1\le p\le\infty$.
\begin{enumerate}
  \item If $u,v\in W^{1,p}(\Omega)\cap L^\infty(\Omega)$, then
        $uv\in W^{1,p}(\Omega)$ and
        \[
          \partial_i(uv)=u\,\partial_i v+v\,\partial_i u.
        \]
  \item If $G\in C^1(\mathbb R)$, $G(0)=0$, and $G'$ is bounded, then
        $G\circ u\in W^{1,p}(\Omega)$ and
        \[
          \partial_i(G\circ u)=(G'\circ u)\partial_i u.
        \]
  \item Let $\Phi\colon\Omega'\to\Omega$ be a $C^1$ bi-Lipschitz
        diffeomorphism whose derivative and inverse derivative are bounded.
        If $u\in W^{1,p}(\Omega)$, then $u\circ\Phi\in W^{1,p}(\Omega')$ and
        \[
          \nabla(u\circ\Phi)
          =D\Phi^{\mathsf T}(\nabla u)\circ\Phi
        \]
        almost everywhere.
  \item If $\eta\in C_c^\infty(\Omega)$ and
        $u\in W^{1,p}_{\mathrm{loc}}(\Omega)$, then
        $\eta u\in W^{1,p}(\Omega)$ and the product rule holds.
\end{enumerate}
\end{noteproposition}
\begin{notedefinition}
For an integer $m\ge1$, a multi-index $\alpha$, and $1\le p\le\infty$,
\[
  W^{m,p}(\Omega)
  =\{u\in L^p(\Omega):D^\alpha u\in L^p(\Omega)
                      \text{ for all }|\alpha|\le m\}.
\]
It is equipped with
\[
  \|u\|_{W^{m,p}}
  =\sum_{|\alpha|\le m}\|D^\alpha u\|_p.
\]
Equivalent finite-dimensional combinations of these terms define the same
norm topology; for $p=2$ this is the Hilbert space
$H^m(\Omega)=W^{m,2}(\Omega)$.
\end{notedefinition}
% Source PDF p. 90 is already marked immediately above; omitting a second
% zero-height marker here prevents an otherwise blank continuation page.
